\subsection{综合练习：从一个数字到完整神经网络}

本附录统一采用三级难度：
\textbf{基础}是回忆或一步手算，
\textbf{标准}是多步复算与形状连接，
\textbf{挑战}是迁移、误差分析或端到端求导。
题目前还保留\textbf{回忆／手算／迁移}的能力类型标签。
可以分次完成，不建议第一次阅读就硬解全部挑战题。

\subsubsection{先做这些题}

\begin{enumerate}[label=\textbf{B1-\arabic*.},leftmargin=2.8em]
\item \textbf{【基础】} \textbf{【手算】} 直线模型 $\widehat y=wx+b$，取 $w=2,b=1$。当 $x=-1,0,2$ 时，预测分别是多少？
\item \textbf{【标准】} \textbf{【手算】} 沿用三点 $(-1,-1),(0,1),(1,3)$，初值为 $\widehat y=x$。求均方误差。
\item \textbf{【标准】} \textbf{【手算】} 对同一三点损失 $L(w,b)=\tfrac23(w-2)^2+(b-1)^2$，从 $(w,b)=(1,0)$ 出发、学习率 $\eta=0.3$。求两项梯度、更新后参数与损失。
\item \textbf{【标准】} \textbf{【手算】} 前向计算 $h=w_1x$、$\widehat y=w_2h$、$L=\tfrac12(\widehat y-y)^2$。若 $x=3,y=8,w_1=2,w_2=1$，求 $h,\widehat y,L$ 以及 $\partial L/\partial w_2$。
\item \textbf{【挑战】} \textbf{【迁移】} 训练误差几乎为零，但验证误差持续变大，说明可能出现什么？至少说出两个改善办法。
\item \textbf{【基础】} \textbf{【手算】} $4\times4$ 图片上使用 $2\times2$ 卷积核，stride 为 1、padding 为 0，输出图像宽高是多少？如果卷积核只有一个输入通道和一个输出通道，权重有几个？
\item \textbf{【基础】} \textbf{【回忆】} 小词表只有四个 \conceptlink{token}{Token}，嵌入维度为 3。嵌入矩阵 $E$ 的形状是多少？一个 Token ID 对应的向量是几维？
\item \textbf{【标准】} \textbf{【手算】} 递推模型 $h_t=x_t+\frac12h_{t-1}$，初始 $h_0=0$，输入为 $(2,0,1)$。计算 $h_1,h_2,h_3$。
\item \textbf{【基础】} \textbf{【手算】} \conceptlink{attention}{Attention} 中 $Q:3\times4$、$K:5\times4$、$V:5\times6$。写出 $QK^\top$、\conceptlink{softmax}{Softmax} 权重及最终输出的形状。
\item \textbf{【标准】} \textbf{【手算】} 一个头输出 $3\times2$ 的矩阵，两个头按特征列拼接后是什么尺寸？若要投影回 4 维，需要什么尺寸的 $W_O$？
\item \textbf{【基础】} \textbf{【手算】} 若某 Token 的真实下一词概率为 $0.25$，\conceptlink{cross-entropy}{交叉熵}单位置损失 $-\log p$ 约为多少？
\item \textbf{【挑战】} \textbf{【迁移】} 给模型加一个私有知识库后，是否能保证它永远不会生成错误事实？解释检索与生成分别可能出什么问题。
\item \textbf{【挑战】} \textbf{【迁移】} 综合数值题（Attention + Softmax + 损失）：
设单 Query 行向量 $q=(1,1)$，
$K=V=\begin{bmatrix}1&0\\0&1\\1&1\end{bmatrix}$。
不使用缩放与掩码。求三个点积分数、Softmax 权重、读取向量；
再让读取行向量 $r$ 右乘一个\textbf{已给定的输出头}：
\[
W_{\rm vocab}=\begin{bmatrix}
2&0&1\\0&1&-1
\end{bmatrix}\in\R^{2\times3}.
\]
求三个 logits $z=rW_{\rm vocab}$、逐行 Softmax 后第二类的概率；
真实标签是第 2 类（从 1 开始计数），求交叉熵
$L=-\log p_2$。解释为什么\textbf{没有输出头矩阵}时，
Attention 权重自身不足以唯一决定候选词概率。
\item \textbf{【挑战】} \textbf{【迁移】}沿用正文统一的 $5\times3$ Embedding 和位置参数。
把输入 $[1,3,4]$ 中第一个 ID 改为 $4$，成为 $[4,3,4]$。
已知 $E[4]=(0.15,-0.2,1.25)$、第一位置向量
$P_1=(0.10,0.05,-0.10)$。
求 Block 的新输入第一行 $X'_1=E[4]+P_1$。
第 2、3 位置的输出是否\textbf{可能}发生变化？
这属于重新训练还是保持权重的输入消融？
\item \textbf{【标准】} \textbf{【手算】} 取不同于正文的
$X=\begin{bmatrix}1&1\\1&-1\end{bmatrix}$。
直接点积求 $XX^\top$，验证原始分数矩阵是否对称；
这与使用独立 $W_Q,W_K$ 的结果有什么不同？
\item \textbf{【标准】} \textbf{【手算】} 把 BPE 的玩具训练语料改为
\texttt{a a b} 与 \texttt{a b}（各为独立序列）。
统计第一轮的相邻对，选最高频对合并；
写出合并后两条序列，并说明是否跨序列计数。
\item \phantomsection\label{exercise:warmup-schedule}\textbf{【标准】} \textbf{【手算】}
给定 $w_0=1$、三步人为固定的梯度 $g_t=2$，
学习率依次为 $(0.01,0.02,0.03)$，按
$w_t=w_{t-1}-\eta_tg_t$ 求 $w_1,w_2,w_3$；
指出这里的 Warmup 改变的是损失函数还是步长。
\item \phantomsection\label{exercise:adamw-decay}\textbf{【标准】} \textbf{【迁移】}
一个参数的 AdamW 步骤给定 $w_0=4$、$\eta=0.1$、
$\lambda=0.05$，偏置校正后的自适应方向为 $0.5$，
忽略 $\varepsilon$。计算更新后的 $w_1$，再令 $\lambda=0$
比较差异；说明权重衰减是否先进入 Adam 的梯度动量。
\end{enumerate}

\noindent\emph{参考答案统一放在\hyperref[app:answers]{附录答案}，不要先翻答案。}

\subsubsection*{给自己画一条完整的路线}
如果你能自己解释下列链条，就掌握了贯穿全书的核心：
\[
\boxed{\text{数据}\to\text{向量}\to\text{函数}\to\text{预测}\to
\text{损失}\to\text{梯度}\to\text{训练}}
\]
\[
\boxed{\text{Token}\to\text{Embedding}\to
\text{Attention/MLP}\to\text{Softmax}\to\text{下一个 Token}}
\]

\begin{conclusion}
先解释问题，再手算一遍，最后用矩阵写短。复杂网络并不是一套全新的算术规则，而是把简单的函数、矩阵乘法、概率和优化方法按任务组合起来。
\end{conclusion}
